\documentclass[../../../include/open-logic-section]{subfiles}

\begin{document}

\olfileid[tr]{sth}{cardinals}{zfc}	
\olsection{$\ZFC$: Bir Dönüm Noktası}

İyi Sıralama'nın eklenmesiyle son teorik dönüm noktasına ulaştık. Artık
$\ZFC$ için gerekli bütün aksiyomlara sahibiz. Ayrıntılı olarak:

\begin{defn}
$\ZFC$ teorisinin aksiyomları Kaplamsallık, Birleşim, Çiftler, Kuvvet
Kümeleri, Sonsuzluk, Temellendirme, İyi Sıralama ve Ayırma ile Yerine Koyma
şemalarının bütün örnekleridir. Başka bir deyişle $\ZFC$, $\ZF$'ye İyi
Sıralama'yı ekler.
\end{defn}

$\ZFC$, \emph{Seçim} ile birlikte \emph{Zermelo--Fraenkel} küme teorisini
belirtir. Bu biraz garip görünebilir; çünkü eklediğimiz aksiyomun adı
``Seçim'' değil ``İyi Sıralama''ydı. Ancak daha sonra Seçim'i formüle
ettiğimizde İyi Sıralama'nın ($\ZF$ altında) Seçim'e denk olduğu ortaya
çıkacaktır (\olref[choice][woproblem]{thmwochoice} sonucuna bakın).
Dolayısıyla hangisini ``temel'' aksiyomumuz olarak alacağımız fark etmez.
Ayrıca ``$\ZFC$'' adı literatürde bütünüyle standarttır.

\end{document}
